\documentclass[11pt]{letter}
\usepackage[margin=1in]{geometry}

\signature{Khaled Eltokhy}
\address{Khaled Eltokhy\\Department of Economics\\The Graduate Center, CUNY}

\begin{document}

\begin{letter}{Editors\\Journal of Economic Behavior \& Organization}

\opening{Dear Editors,}

Please consider my manuscript, ``Speaking in Code: How Surveillance Suppresses Coordinated Dissent,'' for publication in the \textit{Journal of Economic Behavior \& Organization}.

The paper asks how surveillance suppresses collective action in a global game of coordination. Field evidence shows that monitoring chills dissent, but it is difficult to separate direct deterrence from a communication channel in which citizens self-censor before others decide. I isolate that second channel in a controlled language implementation of the Morris--Shin regime-change game. Agents first reproduce the canonical threshold comparative static from prose private signals. Pre-play communication then turns private information into social information. When monitoring is added only to the sender-side communication prompt, leaving the receiver-side decision prompt unchanged and stateless, participation falls. The central result is a belief--action wedge: surveilled messages leave information present but make it less usable for coordination.

The manuscript fits JEBO's interest in economic behavior, organization, and simulation modeling. It combines a standard model of strategic complementarities with controlled language-mediated experiments, and it uses falsification, prompt-isolation, message-content, cross-task, cross-model, and replication-package checks to separate communication-channel effects from generic deterrence or text artifacts. The broader contribution is to show that censorship and surveillance need not be equivalent information-control instruments: censorship removes information from a channel, while surveillance can leave information in circulation but change its common-language form.

\closing{Sincerely,}

\end{letter}

\end{document}
